\documentclass[10pt, letterpaper]{article}

% ------------------------------------------------------
% Packages
% ------------------------------------------------------
\usepackage[margin=1.2cm]{geometry}
\usepackage{titlesec}
\usepackage{array}
\usepackage{enumitem}
\usepackage{hyperref}
\usepackage[dvipsnames]{xcolor}
\usepackage{fontawesome5}
\usepackage{tabularx}

\hypersetup{
    colorlinks=true,
    urlcolor=MidnightBlue
}

\titleformat{\section}{\large\bfseries}{}{0pt}{}[\vspace{2pt}\titlerule]
\titlespacing{\section}{0pt}{0.35cm}{0.15cm}

\renewcommand\labelitemi{$\bullet$}

\setlength{\parskip}{0pt}
\setlength{\parindent}{0pt}

\begin{document}

% ------------------------------------------------------
% Header
% ------------------------------------------------------
\begin{center}
    {\fontsize{22pt}{22pt}\bfseries Mohamed Issaoui}

    \vspace{0.1cm}

    {\small
    \faMapMarker* Sousse, Tunisia \quad
    \faPhone* +216 20 737 821 \quad
    \href{mailto:mohamed.issaoui@ensi-uma.tn}{\faEnvelope[regular] \ mohamed.issaoui@ensi-uma.tn} \quad
    \href{https://www.linkedin.com/in/mohamed-issaoui-/}{\faLinkedinIn \ mohamed-issaoui} \quad
    \href{https://github.com/Medissaoui07}{\faGithub \ Medissaoui07}
    }
\end{center}

\vspace{-0.3cm}

% ------------------------------------------------------
\section{Profile}
Final-year \textbf{Computer Science Engineering Student (AI Track)} at ENSI (national competitive exam rank 56/650), graduating in \textbf{September 2026}.  
Interested in \textbf{deep learning, reinforcement learning, embodied AI, and vision-language navigation for autonomous robots}.  


% ------------------------------------------------------
\section{Education}
\begin{tabularx}{\textwidth}{@{}Xr@{}}
\textbf{Engineering Preparatory Institute of Monastir} – National Competitive Exam Rank: 56/650 & \textit{2021–2023} \\
\textbf{National School of Computer Science (ENSI), Tunisia} – Engineering Degree (AI Track) & \textit{2023–2026} \\
\end{tabularx}

% ------------------------------------------------------
\section{Experience}

\begin{tabularx}{\textwidth}{@{}Xr@{}}
\textbf{AI/Robotics Engineer Intern – OORB, Tunisia} & \textit{Jun 2025 – Aug 2025} \\
\end{tabularx}

\begin{itemize}[leftmargin=0.5cm, itemsep=1pt]
    \item Designed \textbf{quadruped locomotion and manipulation RL environments} with explicit state and reward representations in MuJoCo/Gymnasium.
    \item Trained and deployed \textbf{PPO-based policies} integrated into ROS2 nodes for real-time control.
    \item Built reproducible and scalable experimentation pipelines using \textbf{Docker, Kubeflow, and TensorBoard} on \textbf{Linux and ROS2}.
\end{itemize}

\vspace{0.15cm}

\begin{tabularx}{\textwidth}{@{}Xr@{}}
\textbf{Research Intern (Embodied AI / VLN) – SPIN Lab, Wilfrid Laurier University, Canada} & \textit{Apr 2026 – Jul 2026} \\
\end{tabularx}

\begin{itemize}[leftmargin=0.5cm, itemsep=1pt]
    \item Designed a \textbf{modular Vision-Language Navigation (VLN)} research prototype targeting later deployment on a real \textbf{ROS2} robot.
    \item Built a \textbf{predictive world model} maintaining probabilistic beliefs over unseen spaces, combining \textbf{semantic priors, hypothesis generation, and belief updating} to guide navigation.
    \item Integrated \textbf{VLM-based perception} producing a grounded \textbf{scene-graph memory}, used by an \textbf{LLM planner} for subgoal generation and a stateful navigation policy with frontier scoring.
    \item Experimented with a \textbf{generative anticipation module} to imagine likely future observations under high uncertainty, feeding hypotheses back into the belief state. \textit{Co-authoring a research paper on vision-language navigation (in preparation).}
\end{itemize}

% ------------------------------------------------------
\section{Projects}

\begin{tabularx}{\textwidth}{@{}Xr@{}}
\textbf{Autonomous Mobile Robot – Eurobot Competition} & \textit{2024} \\
\end{tabularx}

\begin{itemize}[leftmargin=0.5cm, itemsep=1pt]
    \item Implemented \textbf{YOLO-based perception} and LiDAR-based SLAM for environment understanding.
    \item Designed \textbf{PID control pipelines} using encoder feedback for reliable navigation.
    \item Deployed complete robotic stacks on \textbf{Raspberry Pi and Jetson Nano} using ROS2.
\end{itemize}

\vspace{0.1cm}

\begin{tabularx}{\textwidth}{@{}Xr@{}}
\textbf{LLM-Guided Robot Learning – ENSI \& OORB} & \textit{2025} \\
\href{https://github.com/PCDorg/RewGen4Robotics}{\faGithub\ GitHub Repo} &
\end{tabularx}

\begin{itemize}[leftmargin=0.5cm, itemsep=1pt]
    \item Developed an \textbf{LLM-guided reinforcement learning framework} inspired by NVIDIA Eureka.
    \item Used language models to \textbf{iteratively reason over environment feedback} and refine reward functions.
    \item Evaluated reward efficiency and behavior quality on MuJoCo robotic tasks with sparse rewards.
\end{itemize}

% ------------------------------------------------------
\section{Technical Skills}
\textbf{Representation \& Reasoning:} Reward modeling, environment abstraction, decision pipelines, world models, scene graphs, belief-state estimation \\
\textbf{Robot Learning:} Reinforcement Learning, imitation learning, LLM-guided reward design \\
\textbf{Deep Learning:} PyTorch, CNNs, Transformers, Vision-Language Models, LLM planners, vision-language navigation \\
\textbf{Robotics:} ROS2, Gazebo, MuJoCo, SLAM, Control Systems, OpenCV, embodied simulation \\
\textbf{Software:} Linux (Ubuntu), Docker, TensorBoard, Kubeflow, Git \\
\textbf{Languages:} Python, C++, C

% ------------------------------------------------------
\section{Languages}
Arabic (Native) \quad French (Fluent) \quad English (Fluent)

% ------------------------------------------------------
\section{Certifications}
Machine Learning Specialization – Coursera \\
NVIDIA Jetson Nano AI Development

\end{document}
